% ---------- Sprache & Kodierung ----------
% ----- KOMMENTAR: nutze fontspec, falls du mit LuaLaTeX oder XeLaTeX kompilierst
\usepackage{fontspec}
\usepackage[german, english]{babel}
\usepackage{microtype}        % optischer Randausgleich, Proportionalität

% =========================================
% SEITENLAYOUT (Lehrbuch-Stil)
% =========================================
\usepackage[
  a4paper,
  twoside, 
  inner=1.5cm,    
  outer=1.5cm,
  top=2.5cm,
  bottom=3cm,
  headheight=14pt,
  includeheadfoot
]{geometry}

\usepackage{setspace}         % Zeilenabstand
% \OnehalfSpacing               % memoir-kompatible Version für 1.5-zeilig
\setlength{\emergencystretch}{3em} % Zeilenumbruch bei Überfüllung erlauben

% ==============================
% Farben & Schrift
% ==============================
\usepackage[dvipsnames]{xcolor}
\definecolor{chaptergray}{gray}{0.20}
\definecolor{linegray}{gray}{0.7}

% ==============================
% Kapitel
% ==============================
\usepackage{titlesec}

% --- Kapitel (ohne Nummer) ---
\titleformat{\chapter}[display]
  {\sffamily\bfseries\LARGE\centering} % Sans-Serif, fett, groß, zentriert
  {} % keine Nummer standardmäßig
  {0pt} % Abstand zwischen Label und Titel (nicht genutzt)
  {} % Vorblock leer
\titlespacing*{\chapter}
  {0pt}   % linker Einzug
  {-10em} % Abstand vor der Überschrift (negativ zieht es nach oben)
  {5ex}   % Abstand nach der Überschrift
% ==============================
% Section
% ==============================
% \titleformat{\section}
%   {\sffamily\bfseries\Large} % Sans-Serif, fett, groß
%   {} % keine Nummer
%   {0pt}
%   {\color{chaptergray}}
% \titlespacing*{\section}{0pt}{2ex}{1ex}

% --- Section nummeriert optional ---
\newcommand{\sectionnumbered}[1]{%
  \refstepcounter{section}%
  \addcontentsline{toc}{section}{\protect\numberline{\thesection}#1}%
  \section*{#1}%
}

% ==============================
% Subsection
% ==============================
\titleformat{\subsection}
  {\sffamily\bfseries\large} % Sans-Serif, fett, etwas kleiner
  {} % keine Nummer
  {0pt}
  {\color{chaptergray}\MakeUppercase} % Großbuchstaben + Farbe
\titlespacing*{\subsection}{0pt}{2ex}{1ex}

% ==============================
% Subsubsection
% ==============================
\titleformat{\subsubsection}
  {\sffamily\bfseries\normalsize} % Sans-Serif, fett
  {} % keine Nummer
  {0pt}
  {\color{chaptergray}}
\titlespacing*{\subsubsection}{0pt}{1.5ex}{0.8ex}

% ==============================
% Kopf- und Fußzeilen für twocolumn
% ==============================
\usepackage{fancyhdr}
\pagestyle{fancy}
\fancyhf{}
\fancyhead[LE]{\small\nouppercase{\leftmark}} % Kapitel links
\fancyhead[RO]{\small\nouppercase{\rightmark}} % Section rechts
\fancyhead[LO,RE]{\thepage}                     % Seitenzahl außen
\renewcommand{\headrulewidth}{0.3pt}
\renewcommand{\footrulewidth}{0pt}

% Kopfzeilen ohne Nummern
\renewcommand{\sectionmark}[1]{\markright{#1}}

% ==============================
% Allgemeine Abstände
% ==============================
\setlength{\parindent}{0em}
\setlength{\parskip}{0.5em}

% \twocolumn  % im main.tex nach \mainmatter aktivieren
\setlength{\columnsep}{0.8cm} % Abstand zwischen Spalten

% optional: optische Linie auf jeder Seite außen
% \usepackage{background}
% \backgroundsetup{position={current page.west}, angle=0, scale=1,
%   color=linegray, opacity=0.3, contents={\rule{0.2pt}{\paperheight}}}

% ---------- Grafiken & Tabellen ----------
\usepackage{graphicx}
% \usepackage{caption}  % wird von subcaption mit kompiliert, daher hier auskommentiert
\usepackage{subcaption}
% \usepackage{float}     % falls [H] benutzt wird
\usepackage{tabularray}
\usepackage{adjustbox}
\usepackage{booktabs}

\renewenvironment{table}[1][Eine Tabelle]{%
  \SetTblrInner{rowsep=5pt}
  \begin{center}
  \begin{adjustbox}{max width=\textwidth}
  \begin{talltblr}[#1]{rows={bg=azure9}, row{1}={bg=azure8}, hlines={white}, vlines={white}}
}{%
  \end{talltblr}
  \end{adjustbox}
  \end{center}
}

% ---------- Sonstige nützliche Pakete ----------
\usepackage{enumitem}
\usepackage[version=4]{mhchem}
\usepackage{tcolorbox}
\usepackage{siunitx}

% Instead of minted, use package listings for code snippets (no external dependencies)

% ---------- Literatur ----------
% Style "numeric" good in general, for physics using AIP-style, i.e "phys"
\usepackage[backend=biber, 
            style=phys, 
            eprint=true,
            sorting=none]{biblatex}
\addbibresource{backmatter/bibliography.bib}
\usepackage{csquotes}

% ---------- Inhaltsverzeichnis ----------
\usepackage{titletoc}
\dottedcontents{section}[5.5em]{}{3em}{1pc}

% ---------- Hyperlinks ----------
\usepackage[
  colorlinks=true,
  linkcolor=blue,
  citecolor=black,
  urlcolor=blue
]{hyperref}
\usepackage{cleveref}   % grouped references

% ---------- ImportAll-Befehl ----------
\usepackage{pgffor}
\newcommand*{\MaxNumOfChapters}{10}
\newcommand*{\MaxNumOfSections}{5}
\newcommand{\importAll}{
    \foreach \c in {1,...,\MaxNumOfChapters}{%
        % Keine Kapitelzählung mehr – volle manuelle Kontrolle
        \foreach \s in {0,1,2,...,\MaxNumOfSections}{%
            \IfFileExists{mainmatter/chapter\c/section\s.tex}{%
                \input{mainmatter/chapter\c/section\s.tex}%
                \clearpage
            }{%
                % Datei existiert nicht → tue nichts
            }%
        }%
    }%
}


% ---------- Kürzel ----------
\newcommand{\red}[1]{\textcolor{red}{#1}}
